\documentclass[11pt]{article}
\usepackage[margin=1in]{geometry}
\usepackage{amsmath,amssymb,amsthm,booktabs,array}
\usepackage[T1]{fontenc}
\usepackage[colorlinks=true,linkcolor=blue,citecolor=blue,urlcolor=blue]{hyperref}
\newtheorem{theorem}{Theorem}
\newtheorem{corollary}[theorem]{Corollary}
\title{Sharp thresholds for Dietz, Flowers and Lockard's\\self-conjugate partition questions:\\extending Flowers and Lockard's construction}
\author{Haobo Ma\\\small ChronoAI Pte Ltd, Singapore
\and Wenlin Zhang\\\small National University of Singapore\\\small ChronoAI Pte Ltd, Singapore}
\date{18 November 2026}
\begin{document}
\maketitle
\begin{abstract}
Using Munagi's residue-class criterion (2008, Theorem~3.1) and extending
Flowers and Lockard's construction (2021, Lemma~4.2, p.~467),
we answer Questions~5.12--5.14 of Dietz, Flowers and Lockard's
\emph{A study of $m$-ary partitions whose conjugates are $q$-ary}
(2026, \S5, p.~404) in their self-conjugate setting: among conjugate
$m$-ary partitions, the first odd exponent admitting a self-conjugate
partition for even $m$ is $m+1$, the first exponent admitting a
non-square self-conjugate partition for any $m$ is $m+1$, and odd
exponents are impossible for odd $m$.
A row-multiplicity identity gives
$N\equiv H+W-1\pmod{(m-1)(q-1)}$ for conjugate $(m,q)$-ary partitions.
For partitions of $m^r$ it yields a positive exponent defect divisible
by $m-1$ whenever the diagram is not rectangular. This proves the exact
nonrectangular threshold and, with the credited constructions, classifies
all exponents admitting a non-square self-conjugate partition.
\end{abstract}

\section{The questions and the congruence}
Dietz, Flowers and Lockard ask three questions about self-conjugate
partitions of powers \cite[Questions~5.12--5.14, p.~404]{DFL}.
The same congruence answers all three in their self-conjugate setting
and proves the rectangle-only range suggested by Flowers and Lockard
\cite[\S4, pp.~466--467]{FL}. The earlier papers construct the examples
at the thresholds using multiplicities and generating operators.
Here a product-modulus congruence supplies matching lower bounds.
For even bases it also explains why a non-square example of an even
power first occurs at exponent $2m$, rather than $m+1$.
The multiplicity step is Munagi's common-residue criterion
\cite[Theorem~3.1, p.~2494]{Munagi}, applied separately to the row and
column moduli. Weighting those multiplicities retains their full product,
even when the moduli have common factors. When they are coprime,
the product congruence also follows by combining the two congruences in
\cite[Lemma~4.1(2)--(3), p.~397]{DFL}.

A partition of a positive integer $N$ is a finite weakly decreasing
sequence $\lambda=(\lambda_1,\ldots,\lambda_H)$ of positive integers with
$\sum_i\lambda_i=N$. Its width is $W=\lambda_1$ and its height is $H$.
Its conjugate has parts
\[
 \lambda'_j=\#\{i:\lambda_i\geq j\},\qquad 1\leq j\leq W.
\]
This reflects the Ferrers diagram in its main diagonal
\cite[\S26.9(i)]{DLMF}. The partition is self-conjugate if
$\lambda'=\lambda$; then $H=W$. A diagram is rectangular if all its rows
have the same length, and square if it is rectangular with $H=W$.
For $m\geq2$, a partition is $m$-ary if every part is a nonnegative
integral power of $m$. A \emph{conjugate $m$-ary partition} (CMP) has both
$\lambda$ and $\lambda'$ $m$-ary. More generally, a conjugate $(m,q)$-ary
partition has $m$-ary rows and $q$-ary columns
\cite[Definitions~1.1--1.2]{DFL}. All bases and exponents below are integers.

\begin{theorem}\label{thm:product}
Let $m,q\geq2$. Every nonempty conjugate $(m,q)$-ary partition of $N$,
with width $W$ and height $H$, satisfies
\begin{equation}\label{eq:product}
 N\equiv H+W-1\pmod{(m-1)(q-1)}.
\end{equation}
For a self-conjugate CMP this gives
\begin{equation}\label{eq:square}
 N\equiv 2W-1\pmod{(m-1)^2}.
\end{equation}
If $m$ is odd, it also gives $N\equiv1\pmod{2(m-1)}$.
\end{theorem}
\begin{proof}
Let $S$ be the finite set of distinct row lengths, $a_w$ the multiplicity
of $w\in S$, and $c_j=\#\{i:\lambda_i\geq j\}$ for $1\leq j\leq W$,
with $c_{W+1}=0$. Then $a_w=c_w-c_{w+1}$,
$H=\sum_{w\in S}a_w$ and $N=\sum_{w\in S}a_ww$.
Every $w\in S$ is a power of $m$, so $m-1\mid w-1$.
If $w<W$, both $c_w$ and $c_{w+1}$ are positive powers of $q$;
therefore $q-1\mid a_w$. At $w=W$, $a_W=c_W$ is a power of $q$,
so $q-1\mid a_W-1$. Put $\varepsilon_W=1$ and
$\varepsilon_w=0$ for $w<W$. Each summand in
\[
 \sum_{w\in S}(a_w-\varepsilon_w)(w-1)
 =N-H-W+1
\]
is divisible by $(m-1)(q-1)$, proving \eqref{eq:product}.
Self-conjugacy gives $H=W$ and hence \eqref{eq:square}.
For odd $m$, put $d=m-1$. Since $d$ is even, $2d\mid d^2$;
also $d\mid W-1$ because $W$ is a power of $m$.
Thus $N-1=(N-2W+1)+2(W-1)$ is divisible by $2d$.
\end{proof}

\section{The first nonrectangular partition of a power}
\begin{corollary}\label{cor:defect}
Let a nonrectangular CMP of $m^r$, $m\geq2$, have width $m^a$
and height $m^b$. Then
\[
 1\leq a+b-r\leq r-2,\qquad m-1\mid a+b-r.
\]
In particular $r\geq m+1$. If it is self-conjugate, $a=b$ and
\begin{equation}\label{eq:defect}
 1\leq 2a-r\leq r-2,\qquad m-1\mid 2a-r.
\end{equation}
\end{corollary}
\begin{proof}
Write $d=m-1$. The binomial theorem gives
$m^j=(1+d)^j\equiv1+jd\pmod{d^2}$ for every $j\geq0$.
Theorem~\ref{thm:product} implies $d^2\mid d(r-a-b)$, hence
$d\mid a+b-r$, including $d=1$.
Nonrectangularity gives $m^r<m^{a+b}$.
The first row and column give $m^r\geq m^a+m^b-1$.
A nonrectangular diagram has $a,b\geq1$, so
$m^r>\max(m^a,m^b)$ and $a,b\leq r-1$.
These bounds give the two inequalities. A positive multiple of $d$
is at least $d$, so $d\leq r-2$.
\end{proof}

\begin{corollary}\label{cor:rectangles}
A nonrectangular CMP of $m^r$, for $m\geq2$ and $r\geq0$, exists
exactly when $r\geq m+1$. If $C_m(N)$ counts the CMPs of $N$, then
\[
 C_m(m^r)=r+1\quad(0\leq r\leq m),\qquad
 C_m(m^r)>r+1\quad(r\geq m+1).
\]
The smallest non-square self-conjugate CMP of a power of $m$ has size
$m^{m+1}$.
\end{corollary}
\begin{proof}
Corollary~\ref{cor:defect} gives the lower bound.
For attainment, Flowers and Lockard's partition
\cite[Lemma~4.2, p.~467]{FL} has one row of length $m^m$, and
$m^j-m^{j-1}$ rows of length $m^{m-j}$ for $1\leq j\leq m$.
The number of rows of length at least $m^{m-j}$ is $m^j$.
Thus conjugation gives the same parts and multiplicities. Its size is
\[
 m^m+\sum_{j=1}^m(m^j-m^{j-1})m^{m-j}
 =m^m+m(m-1)m^{m-1}=m^{m+1}.
\]
It is non-square and self-conjugate.
Multiplying each row length by $m^t$ repeats each conjugate part
$m^t$ times, preserving the CMP property and nonrectangularity;
the size becomes $m^{m+1+t}$.
The rectangular CMPs of $m^r$ are exactly $m^{r-j}$ rows of length
$m^j$, $0\leq j\leq r$. These $r+1$ rectangles are all distinct.
Finally, a self-conjugate rectangle is a square.
\end{proof}

\section{Self-conjugate powers and sharp parity thresholds}
\begin{corollary}\label{cor:negative}
For every odd $m\geq3$ and every $s\geq0$, no self-conjugate CMP of
$m^{2s+1}$ exists. A self-conjugate CMP of $m^r$, $r\geq0$, exists
if and only if $r$ is even.
\end{corollary}
\begin{proof}
Set $d=m-1$. Since $m+1$ is even, $m^2\equiv1\pmod{2d}$, so
$m^{2s+1}\equiv m\pmod{2d}$. Theorem~\ref{thm:product} would instead
give $m^{2s+1}\equiv1\pmod{2d}$, requiring $2d\mid d$, which is
impossible for $d>0$. If $r=2k$, the square with $m^k$ rows of
length $m^k$ supplies existence, including $(1)$ for $k=0$.
This square construction is recalled in
\cite[Discussion~5.11, p.~404]{DFL}.
\end{proof}
For the subfamily $m\equiv3\pmod4$, nonexistence also follows from the
earlier odd-parts restriction implicit in Munagi's generating function
\cite[Theorem~5.2, p.~2500]{Munagi}: every exponent in
$1+\sum_{n\geq0}x^{(2n+1)^2}/\prod_{j=1}^{2n+1}(1-x^{4j})$
is zero or $1\pmod4$. Somos records the corresponding size restriction
in his dated OEIS comment \cite{OEIS}; Andrews subsequently studies
the same class \cite[\S3, Lemma~3.1]{Andrews}.
The modulus $2(m-1)$ in Theorem~\ref{thm:product} excludes all odd bases,
including $m\equiv1\pmod4$.

\begin{corollary}\label{cor:classification}
A non-square self-conjugate CMP of $m^r$, $m\geq2$, $r\geq0$, exists
exactly in these cases:
\[
\begin{array}{c|c}
\text{base and exponent parity}&\text{necessary and sufficient range}\\\hline
m\text{ odd},\ r\text{ even}&r\geq m+1\\
m\text{ even},\ r\text{ odd}&r\geq m+1\\
m\text{ even},\ r\text{ even}&r\geq2m.
\end{array}
\]
For odd $m$ and odd $r$ none exists. Allowing squares, a self-conjugate
CMP exists precisely when $r$ is even or when $m$ is even and $r\geq m+1$.
\end{corollary}
\begin{proof}
Corollary~\ref{cor:defect} gives $r\geq m+1$, and
Corollary~\ref{cor:negative} excludes odd $m,r$.
If $m,r$ are both even, $2a-r$ in \eqref{eq:defect} is even and
divisible by the odd integer $m-1$. Hence $2a-r\geq2(m-1)$ and
$r\geq2m$.

Corollary~\ref{cor:rectangles} attains $m+1$. For even $m$, Dietz,
Flowers and Lockard construct non-square self-conjugate CMPs of
$m^{2m}$ \cite[Corollaries~5.7--5.10, pp.~403--404]{DFL}; a direct
description follows. Replacing each cell by an $m\times m$ square
preserves self-conjugacy, non-square shape and the CMP property,
and increases $r$ by two \cite[Theorem~5.4, p.~469]{FL}.
These starting exponents attain every range in the table.
Squares supply every even exponent, including zero.
\end{proof}

Here is the construction from \cite[Corollary~5.9, p.~404]{DFL}
in diagram terms. For a self-conjugate CMP of width $W$, let $U$
replace one largest row $W$ by $mW$ and append $(m-1)W$ rows of
length one. This adds equal arms to the first row and column,
preserving self-conjugacy and the CMP property, and takes
$(W,N)$ to $(mW,N+2(m-1)W)$.
Let $T$ replace each cell by an $m\times m$ square, taking
$(W,N)$ to $(mW,m^2N)$. Composition acts from right to left.
Starting with the partition $(1)$, set
$\lambda_j=(U^2T^2)^jU^2T(1)$ for $j\geq0$.
Its width and size are
\[
 W_j=m^{4j+3},\qquad
 N_j=2(j+1)m^{4j+3}+m^{4j+2}-2(j+1)m^{4j+1}.
\]
The case $j=0$ follows by applying $T$ and then $U$ twice.
The recurrences $W_{j+1}=m^4W_j$ and
$N_{j+1}=m^4N_j+2(m^2-1)m^2W_j$ give the formulas by induction.
For even $m$, set $j=(m-2)/2$. The last two terms cancel and
$N_j=m^{2m}$. Each diagram is non-square because its last $U$
creates distinct largest and smallest row lengths.

\section{A density consequence for mixed bases}
The general density question in \cite[Question~4.9, p.~400]{DFL}
also has a short partial answer from the size bounds for its operators.
\begin{corollary}\label{cor:density}
Let $A_{m,q}(X)$ count the positive integers at most $X$ having a
conjugate $(m,q)$-ary partition. If
$m^{-1/2}+q^{-1/2}<1$, then $A_{m,q}(X)/X\to0$ as $X\to\infty$.
The same conclusion holds for $(m,q)=(4,4)$.
\end{corollary}
\begin{proof}
By \cite[Theorem~2.7 and Corollary~2.9]{DFL}, every such partition
has a unique generating word with $a$ migration and $b$ raise operators;
its width and height are $m^a$ and $q^b$.
A diagram of size at most $X$ has $m^a,q^b\leq X$.
Choose $0<\sigma<1/2$ with $m^{-\sigma}+q^{-\sigma}<1$.
Since $1\leq X^{2\sigma}m^{-\sigma a}q^{-\sigma b}$,
the number of eligible generating words is at most
\[
 X^{2\sigma}\sum_{a,b\geq0}\binom{a+b}{a}
 m^{-\sigma a}q^{-\sigma b}
 =\frac{X^{2\sigma}}{1-m^{-\sigma}-q^{-\sigma}}=o(X).
\]
This also bounds $A_{m,q}(X)$.
For $m=q=4$, put $L=\lfloor\log_4X\rfloor$; then $a,b\leq L$.
Pascal's identity, summed twice, gives
\[
 A_{4,4}(X)\leq\sum_{a,b=0}^L\binom{a+b}{a}
 =\binom{2L+2}{L+1}-1
 \leq\frac{4^{L+1}}{\sqrt{L+2}}.
\]
For the last bound use
$\binom{2n}{n}/4^n=\prod_{j=1}^n(1-1/(2j))
\leq\exp(-\tfrac12\sum_{j=1}^n1/j)\leq(n+1)^{-1/2}$.
Since $X\geq4^L$, division by $X$ proves the claim.
\end{proof}

\section{Source coverage and mathematical verification}
Question~5.13 says ``smallest non-square CMP'' without repeating
``self-conjugate''. In its section's self-conjugate setting the answer
is yes. Without that setting a non-square rectangle such as $(m)$
already partitions $m$, so the unrestricted reading has answer no.
Corollary~\ref{cor:rectangles} also establishes minimality among all
\emph{nonrectangular} CMPs, without assuming self-conjugacy.
\begin{center}
\small
\begin{tabular}{@{}p{0.32\linewidth}p{0.61\linewidth}@{}}
\toprule
Source item & Coverage\\\midrule
\cite{DFL}, Question 5.12, p.~404 & Answered: yes. The first odd
exponent for even $m$ is $m+1$ (Corollary~\ref{cor:classification}).\\[3pt]
\cite{DFL}, Question 5.13, p.~404 & Answered: yes in the self-conjugate
setting; first non-square exponent $m+1$ (Corollary~\ref{cor:rectangles}).\\[3pt]
\cite{DFL}, Question 5.14, p.~404 & Answered: no
(Corollary~\ref{cor:negative}).\\[3pt]
\cite{DFL}, Discussion 5.11; Corollaries 5.7--5.10 & Prior constructions
cited: Dietz--Flowers--Lockard, 2026; Corollary 5.9 reproved above.
Sharp parity thresholds: Corollary~\ref{cor:classification}; the full
families and their counts are not covered.\\[3pt]
\cite{DFL}, Lemma 4.1, pp.~397--398 & Complemented by the product-modulus
congruence of Theorem~\ref{thm:product}.\\[3pt]
\cite{FL}, Lemma 4.2, p.~467; Theorem 5.4, p.~469 & Reproved; prior
proof Flowers--Lockard, 2021. Constructions used in
Corollaries~\ref{cor:rectangles} and~\ref{cor:classification}.\\[3pt]
\cite{FL}, rectangle-only observation, pp.~466--467 & Proved:
$C_m(m^r)=r+1$ for $0\leq r\leq m$ (Corollary~\ref{cor:rectangles}).\\[3pt]
\cite{Munagi}, Theorems 3.1 and 5.2 & Prior criterion and generating
function cited: Munagi, 2008. The latter already gives the
$m\equiv3\pmod4$ subcase of the odd-power obstruction.\\[3pt]
\cite{DFL}, Question 4.9, p.~400 & Partially answered: density zero
under Corollary~\ref{cor:density}'s hypotheses. Other pairs and exact
existence/counting are not covered.\\[3pt]
\cite{DFL}, Example 4.7, p.~399 & Not covered: the listed unresolved
power-uniqueness questions.\\\bottomrule
\end{tabular}
\end{center}
The other generating-operator, polynomial and counting results in the
two source papers are outside the scope of this note.

The congruences, threshold bounds and existence classifications have
written proofs above, with the earlier constructions credited separately.
The Lean declaration
\begin{quote}\small\ttfamily
D5.S3.Combinatorics.Partitions.\\
OddPowerSelfConjugateCMP.result
\end{quote}
in the \href{https://github.com/the-omega-institute/trureturing/blob/8acf7efa41af7fc9ddbbd39571e7ef73462f4509/D5/S3/Combinatorics/Partitions/OddPowerSelfConjugateCMP.lean}{pinned nonexistence source}
proves precisely the first statement of Corollary~\ref{cor:negative}:
for every odd $m\geq2$ and $s\geq0$, a finite Young diagram of size
$m^{2s+1}$ with $m$-ary rows and columns is not equal to its transpose.
The mixed-base extension, minimality results, existence classifications
and density bound are covered by the written arguments.
In the \href{https://github.com/the-omega-institute/trureturing/blob/8acf7efa41af7fc9ddbbd39571e7ef73462f4509/D5/S1/Words/Compositions/StrictFirstSumsGapfreeOddParts.lean}{pinned structural helper source},
\texttt{rowLen\_zero} and \texttt{height\_bounds} give zero rows beyond
the height and bounds for positive column lengths. They support the
formal nonexistence proof, not the added threshold statements.
The accompanying \texttt{verify.py} checks exact congruences and
transposes over bounded finite sets and checks the constructions using
compressed multiplicities. Its output specifies all finite ranges;
these checks supplement the written proofs.

\begin{thebibliography}{9}
\bibitem{DFL} G.~D. Dietz, T.~B. Flowers and S.~R. Lockard,
\emph{A study of $m$-ary partitions whose conjugates are $q$-ary},
J. Algebra Comb. Discrete Struct. Appl. \textbf{13}(3) (2026), 385--407.
\href{https://doi.org/10.13069/jacodesmath.v13i3.393}{doi:10.13069/jacodesmath.v13i3.393}.
\bibitem{FL} T.~B. Flowers and S.~R. Lockard,
\emph{Conjugate $m$-ary partitions}, Australas. J. Combin.
\textbf{81}(3) (2021), 458--473.
\url{https://ajc.maths.uq.edu.au/pdf/81/ajc_v81_p458.pdf}.
\bibitem{Munagi} A.~O. Munagi, \emph{Pairing conjugate partitions by
residue classes}, Discrete Math. \textbf{308}(12) (2008), 2492--2501.
\href{https://doi.org/10.1016/j.disc.2007.05.022}{doi:10.1016/j.disc.2007.05.022}.
\bibitem{OEIS} OEIS Foundation Inc., \emph{A067357: Number of
self-conjugate partitions of $4n+1$ into odd parts},
comment by M.~Somos dated 25 July 2015.
\url{https://oeis.org/A067357}. Accessed 18 November 2026.
\bibitem{Andrews} G.~E. Andrews, \emph{The Bhargava-Adiga Summation
and Partitions}, J. Indian Math. Soc. \textbf{84}(3--4) (2017), 151--160.
\href{https://doi.org/10.18311/jims/2017/15836}{doi:10.18311/jims/2017/15836}.
\bibitem{DLMF} NIST Digital Library of Mathematical Functions,
\S26.9(i), \emph{Integer partitions: restricted number and part size---Definitions}.
\url{https://dlmf.nist.gov/26.9#i}. Accessed 18 November 2026.
\end{thebibliography}
\medskip
\noindent\small \copyright\ 2026 Haobo Ma and Wenlin Zhang.
This work is licensed under the Creative Commons Attribution 4.0
International license (CC BY 4.0).
\end{document}
